% ===== 原文 PDF pp.35--37（印刷 pp.34--36）=====
% 译者 tr-03
% [原文 p.34]
相反，存在多个（但等价）的拉格朗日量，它们带有不同的基本变量组，彼此之间通过电磁对偶的
超对称版本相联系。在 $\cM_{\mathrm{CB}}$ 上的一般点附近，它们都是低能物理同样好的局域描述。
整体上，CB 有效场论由这样的局域开片拼成，粘合映射包含导致单值性的非平凡对偶变换。这样的
对偶单值性必然要求 CB 上存在奇点，而这正是 CB 具有有趣动力学的原因。在本小节中，我们将
陈述电磁对偶的一些基本要点。

让我们从回顾 4 维纯麦克斯韦理论中的电磁对偶如何运作开始。考虑欧几里得号差下麦克斯韦理论的
配分函数
\begin{equation}
  \int[\cD A_\mu]\exp\int d^4x\left(
  -\frac{1}{4g^2}F_{\mu\nu}F^{\mu\nu}
  + \frac{i\theta}{32\pi^2}\epsilon_{\mu\nu\rho\sigma}F^{\mu\nu}F^{\rho\sigma}\right).
  \tag{3.11}
\end{equation}
对偶变换是通过把积分变量从规范场 $A$ 换成场强 $F$ 来实现的。然而，此时我们其实正在丢失
信息：麦克斯韦理论并不只是关于 2-形式的理论，而是关于局域恰当 2-形式的理论。这一点通常
编码在 Bianchi 恒等式 $dF=0$ 之中。为了在新的对偶表述（对场强空间做路径积分）中恢复这一
信息，我们引入一个新的拉格朗日乘子 1-形式 $A^D$，并写出
\begin{equation}
  \resizebox{0.82\textwidth}{!}{$\displaystyle
  \int[\cD F_{\mu\nu}]\,[\cD A^D_\lambda]\exp\int d^4x\left(
  -\frac{1}{4g^2}F_{\mu\nu}F^{\mu\nu}
  + \frac{i\theta}{32\pi^2}\epsilon^{\mu\nu\rho\sigma}F_{\mu\nu}F_{\rho\sigma}
  + \frac{i}{8\pi}\epsilon^{\mu\nu\rho\sigma}\partial_\mu A^D_\nu F_{\rho\sigma}\right)\,.$}
  \tag{3.12}
\end{equation}
为过渡到对偶描述，我们现在必须对 $F$ 做积分。由于这是高斯型理论，我们可以在经典层面完成
它，即借助 $F$ 的运动方程消去 $F$。如果我们先把拉格朗日量用自对偶与反自对偶场强
$F^\pm_{\mu\nu} := \frac{1}{2}\left(F_{\mu\nu}\pm\frac{1}{2}\epsilon_{\mu\nu\rho\sigma}F^{\rho\sigma}\right)$
改写（$A^D$ 的场强类似），这一操作会更加透明：
\begin{equation}
  \frac{i}{8\pi}\int d^4x\left(\bar\tau(F^+)^2-\tau(F^-)^2\right)
  + \frac{i}{4\pi}\int d^4x\left(F^+_D F^+ - F^-_D F^-\right).
  \tag{3.13}
\end{equation}
请读者验证：把 $F^\pm$ 替换为其运动方程后，即得到如下路径积分
\begin{equation}
  \resizebox{0.82\textwidth}{!}{$\displaystyle
  \int[\cD A^D_\mu]\exp\left(-\frac{i}{8\pi}\int d^4x
  \left[\frac{-1}{\bar\tau}\left(F^+_D\right)^2-\frac{-1}{\tau}\left(F^-_D\right)^2\right]\right)\,.$}
  \tag{3.14}
\end{equation}

% [Tutor Note: 对比 (3.11) 与 (3.12) 的指标位置可见原文的一处记号切换：(3.11) 中写成
% $\epsilon_{\mu\nu\rho\sigma}F^{\mu\nu}F^{\rho\sigma}$（$\epsilon$ 全下标、$F$ 全上标），
% 而 (3.12) 中写成 $\epsilon^{\mu\nu\rho\sigma}F_{\mu\nu}F_{\rho\sigma}$（全上标/全下标）。
% 两者只差整体符号约定，物理内容相同；此处按 800 dpi 裁图逐字符核对后原样复刻，不作统一。]

% [原文 p.35]
它确实描述了同一个麦克斯韦理论，但用的是对偶磁变量 $A_D$ 与对偶耦合
$\tau_D := -\frac{1}{\tau}$。

比较 (3.14) 与 (3.11)，我们得出结论：麦克斯韦理论在如下变换下不变\footnote{顺带一提，我们
想强调：若 4 维时空是弯曲的，则麦克斯韦理论在 $\SL(2,\Z)$ 的某个有限子群下可能存在某种反常
[152,153]，该子群在 $\tau$ 的特殊取值处成为理论的对称性，文献 [154,155] 称之为对偶反常。}：
\begin{equation}
  A\to A_D\,,\qquad \tau\to-\frac{1}{\tau}\,.
  \tag{3.15}
\end{equation}
此外，注意到麦克斯韦理论在 $\theta\to\theta+2\pi n$（$n\in\Z$）的平移下也不变，整个对偶群
被增强为 $\SL(2,\Z)$。

\begin{exercisebox}{习题 3.1}
  把 $F^\pm$ 的运动方程代入式 (3.13)，并证明所得理论就是式 (3.14)。
\end{exercisebox}

% [Tutor Note: 习题 3.1 的解法要点（给读者提示，不进入正文）：把 (3.13) 对 $F^\pm$ 变分，
% 即得 $F^+_D = -\bar\tau^{-1}F^+$、$F^-_D=-\tau^{-1}F^-$（这正是 (3.14) 指数中两个系数的来源）；
% 代回后交叉项给出 $\bar\tau(F^+)^2$ 与 $\tau(F^-)^2$ 的系数 $\pm\frac{i}{8\pi}\cdot(-1)$，
% 与 (3.14) 中的 $\frac{-1}{\bar\tau}$、$\frac{-1}{\tau}$ 对上，符号亦一致。]

到目前为止，我们关注的都是纯麦克斯韦理论中的电磁对偶。在具有 $\cN=2$ 超对称时，人们自然
会预期：与规范场 $A_i$ 处在同一超多重态中的标量场 $a_i$ 也应享有同样的对偶性质。为此，让
我们回到 CB 上的有效理论。关键的洞见是：(3.4) 与 (3.6) 所定义的
$\Im\left(\tau^{\mathrm{IR}}_{ij}\right)$ 既是调和的又是正定的，因此除非它是常数（此时理论是
自由的），否则不可能在整个模空间 $\cM_{\mathrm{CB}}$ 上整体定义。尽管如此，
$\Im\left(\tau^{\mathrm{IR}}_{ij}\right)$ 是局域良定义的，例如在 CB 上的半经典（弱耦合）区域。
把它延拓到整个 CB 的障碍来自真正的量子效应：这些效应修改了经典模空间的结构，而后者先验地
可以由 $a_i$ 整体描述。为简单起见，本节余下部分将只讨论秩一的情形，因此下标 $ij$ 可以省略。

若我们定义另一个坐标
\begin{equation}
  a^D = \frac{\partial\cF}{\partial a}\,,
  \tag{3.16}
\end{equation}
则 CB 上的度规可写成
\begin{equation}
  ds^2 = \Im\,da_D\,d\bar a = -\frac{i}{2}\left(da_D\,d\bar a - da\,d\bar a_D\right).
  \tag{3.17}
\end{equation}
这里 $a$ 与 $a_D$ 是 $u\equiv\Tr(\Phi^2)$ 在 $\cM_{\mathrm{CB}}$ 上的（多值）全纯函数。纯
麦克斯韦理论中 $\SL(2,\Z)$ 对偶的 $\cN=2$ 超对称推广作用在 $(a,a_D)$ 上为
\begin{equation}
  \begin{pmatrix} a^D\\[2pt] a\end{pmatrix}\to
  \begin{pmatrix} f & g\\[2pt] h & l\end{pmatrix}
  \begin{pmatrix} a^D\\[2pt] a\end{pmatrix},
  \tag{3.18}
\end{equation}

% [Tutor Note: 原文在 $a^D$ 的指标写法上并不统一：(3.16)、(3.18)、(3.19)、(3.22)、(3.23) 以及
% p.38 的正文用上标 $a^D$，而 (3.17) 与 p.38 的图 5 图注用下标 $a_D$，p.40 正文两者都出现。
% 已用 800--900 dpi 裁图分别复核 (3.16)/(3.17)/(3.18)/(3.19) 四处，确认不是文字层抽取错误，
% 而是原文自身的记号不一致；本译本按各处原文逐字复刻，读者应把它们理解为同一个对偶坐标。]

% [原文 p.36]
其中 $f,g,h,l\in\Z$，且矩阵的行列式等于一。上述变换显然保持 CB 度规 (3.17) 不变，\footnote{CB
上的特殊 Kähler 度规实际上在更大的 $\SL(2,\R)$ 群（高秩情形为 $\Sp(2r,\R)$）下不变；由于
电磁荷的量子化，该群被约化为 $\SL(2,\Z)$（一般情形为 $\Sp(2r,\Z)$）。}并且它作用在有效耦合
$\tau^{\mathrm{IR}}$ 上，
\begin{equation}
  \tau^{\mathrm{IR}} = \frac{\partial a^D}{\partial a}\,,
  \tag{3.19}
\end{equation}
的方式为
\begin{equation}
  \tau^{\mathrm{IR}}\to\frac{f\tau^{\mathrm{IR}}+g}{h\tau^{\mathrm{IR}}+l}\,,
  \tag{3.20}
\end{equation}
而这正是 $\SL(2,\Z)$ 对偶作用在麦克斯韦耦合上的方式。

粒子态在上述对偶下也会变换，例如电子、磁单极子与双子。在这些对象中，我们可以辨认出满足
某些 BPS 条件 [158] 的 BPS 粒子 [156,157]。的确，回顾式 (2.12) 中中心荷 $\cZ$ 的定义，对任何
质量为 $M$ 的 $\cN=2$ 粒子，由超对称代数可得
\begin{equation}
  M\ge|\cZ|\,,
  \tag{3.21}
\end{equation}
而对 BPS 粒子要求该不等式取等号。注意，这类达到 BPS 饱和的态受 $\cN=2$ 超对称保护，不会因
微扰或非微扰修正而偏离该界限。\footnote{更准确地说，BPS 粒子在 CB 上的一般点是（因 BPS
界而）稳定的。在某些实余维一的位置上存在边际稳定墙，BPS 粒子在其上可能衰变，从而导致 BPS
谱的非平凡跳变，即所谓的穿墙现象。关于 $\SU(2)$ 规范理论中此类现象的描述，参见 [20,21]。}
在 4d $\cN=2$ 理论中，中心荷在半经典（弱耦合）区域由下式确定
\begin{equation}
  \cZ = n\,a + m\,a^D + \sum_A f_A\,\mu_A\,,
  \tag{3.22}
\end{equation}
其中现在 $a$ 与 $a^D$ 是对应（对偶）$\cN=2$ 矢量多重态的最下分量，$\mu_A$ 表示味对称性的
背景质量参数。最后，$n$、$m$ 与 $f_A$ 分别是 BPS 粒子的电、磁与味荷。

在数学上，在 $\cM_{\mathrm{CB}}$ 的不同开片中，二重态
\begin{equation}
  \left(a,\,a^D\right),
  \tag{3.23}
\end{equation}
定义了\textbf{特殊坐标}；正如我们后面将展示的，它们与 $\cM_{\mathrm{CB}}$ 上的特殊几何紧密
相连。

% [Tutor Note: (3.22) 中 $n$ 与 $m$ 是整数（电、磁荷量子数），$f_A$ 是味荷，三者均为无量纲；
% $a$、$a^D$ 与 $\mu_A$ 的量纲都是质量一次，故 $\cZ$ 与 $M$ 同量纲，这正是 (3.21) 的
% $M\ge|\cZ|$ 能成立的前提。求和 $\sum_A$ 遍历所有味对称性 $A$。]
